2\centerline{\bf \Huge Неравенства Hard}

\begin{flushright}
        \scriptsize{Эпизод 23. Слёзы}
\end{flushright}

\problem{1}
Даны числа $a, b, c$, не меньше 1. Докажите, что

$$
\frac{a+b+c}{4} \geqslant \frac{\sqrt{a b-1}}{b+c}+\frac{\sqrt{b c-1}}{c+a}+\frac{\sqrt{c a-1}}{a+b}
$$

\problem{2}
Даны положительные числа $x_1, x_2, \ldots, x_n$, где $n \geqslant 2$. Докажите, что

$$
\frac{1+x_1^2}{1+x_1 x_2}+\frac{1+x_2^2}{1+x_2 x_3}+\ldots+\frac{1+x_{n-1}^2}{1+x_{n-1} x_n}+\frac{1+x_n^2}{1+x_n x_1} \geqslant n .
$$

\problem{3}
Сумма положительных чисел $a, b, c, d$ равна 3. Докажите неравенство:

$$
\frac{1}{a^2}+\frac{1}{b^2}+\frac{1}{c^2}+\frac{1}{d^2} \leqslant \frac{1}{a^2 b^2 c^2 d^2}
$$

\problem{4}
Найдите наибольшее вещественное $a$ такое, что для всех натуральных $n \geqslant 1$ и всех вещественных чисел $0=x_0<x_1<x_2<\ldots<x_n$ выполнено неравенство

$$
\frac{1}{x_1-x_0}+\frac{1}{x_2-x_1}+\ldots+\frac{1}{x_n-x_{n-1}} \geqslant a\left(\frac{2}{x_1}+\frac{3}{x_2}+\ldots+\frac{n+1}{x_n}\right)
$$

\problem{5}
Дано число $a \in(0,1)$. Положительные числа $x_0, x_1, x_2, \ldots, x_n$ удовлетворяют условиям

$$
x_0+x_1+\ldots+x_n=n+a
$$

и

$$
\frac{1}{x_0}+\frac{1}{x_1}+\ldots+\frac{1}{x_n}=n+\frac{1}{a}
$$


Найдите наименьшее возможное значение выражения

$$
x_0^2+x_1^2+\ldots+x_n^2
$$

